\section{ LÝ THUYẾT SƠ LƯỢC}

% \subsection{Kiến trúc RISC}
% \textbf{Đặc điểm chính của kiến trúc RISC bao gồm:}
% \begin{itemize}
%     \item \textbf{Tập lệnh ít và đồng đều:} Mỗi lệnh có độ dài cố định (thường 32-bit), giúp giải mã lệnh nhanh hơn và đơn giản hóa phần cứng giải mã.
%     \item \textbf{Chế độ địa chỉ đơn giản:} Chỉ hỗ trợ Load/Store để truy cập bộ nhớ; các phép toán được thực thi trực tiếp trên thanh ghi.
%     \item \textbf{Nhiều thanh ghi đa năng:} Giảm thiểu số lần truy cập bộ nhớ, tăng tốc độ xử lý.
%     \item \textbf{Pipeline hiệu quả:} Cấu trúc lệnh đều đặn, cố định tạo điều kiện thuận lợi cho xử lý đường ống (pipelining).
% \end{itemize}

% \textbf{Ưu điểm:}
% \begin{itemize}
%     \item Tốc độ thực thi lệnh cao do mỗi lệnh hoàn thành trong một chu kỳ đồng hồ.
%     \item Thiết kế mạch đơn giản hơn, dễ kiểm tra và tối ưu hóa.
%     \item Tiêu thụ năng lượng thấp hơn so với CISC, phù hợp với thiết bị di động và nhúng.
% \end{itemize}

% \textbf{Nhược điểm:}
% \begin{itemize}
%     \item Chương trình máy thường dài hơn do mỗi tác vụ phức tạp cần được phân rã thành nhiều lệnh đơn giản.
%     \item Yêu cầu nhiều thanh ghi hơn, làm tăng diện tích chip.
% \end{itemize}

% \begin{table}[H]
% \centering
% \renewcommand{\arraystretch}{1.2}
% \small
% \begin{tabular}{|l|l|l|}
% \hline
% \textbf{Tiêu chí} & \textbf{RISC} & \textbf{CISC} \\ \hline
% Tập lệnh & Ít, đơn giản, cố định & Nhiều, phức tạp, biến thiên \\ \hline
% Độ dài lệnh & Cố định & Thay đổi \\ \hline
% Chu kỳ thực thi & 1 chu kỳ/lệnh & Nhiều chu kỳ/lệnh \\ \hline
% Bộ giải mã & Đơn giản & Phức tạp \\ \hline
% Thanh ghi & Nhiều & Ít hơn \\ \hline
% Pipeline & Dễ triển khai & Khó triển khai \\ \hline
% \end{tabular}
% \caption{So sánh kiến trúc RISC và CISC}
% \end{table}


\subsection{Cấu trúc RISC CPU}
Một RISC CPU đơn giản bao gồm các khối chức năng chính sau:
\begin{itemize}
    \item \textbf{ALU (Arithmetic Logic Unit):} Thực thi các phép toán số học (cộng, trừ) và logic (AND, XOR, NOT) trên dữ liệu đầu vào.
    \item \textbf{Accumulator:} Thanh ghi lưu trữ tạm thời toán hạng và kết quả trung gian trong quá trình xử lý.
    \item \textbf{Controller:} Khối điều khiển giải mã opcode và điều khiển các tín hiệu đến các khối khác.
    \item \textbf{Address Mux: } Khối đa hợp lựa chọn giữa địa chỉ của PC hoặc IR theo sự điều khiển của khối Controller.
    \item \textbf{Program Counter (PC):} Thanh ghi lưu địa chỉ của lệnh tiếp theo cần thực thi; tự động tăng sau mỗi chu kỳ lệnh.
    \item \textbf{Instruction Register (IR):} Lưu mã lệnh vừa được nạp từ bộ nhớ, phục vụ cho quá trình giải mã.
    \item \textbf{Memory:} Bộ nhớ lưu trữ chương trình (instruction memory) và dữ liệu (data memory).
\end{itemize}

\textbf{Chu trình hoạt động Fetch - Decode - Execute (FDE):}
Mặc dù khối điều khiển được thiết kế dưới để hoạt động trong 8 pha (8-phase FSM), 
toàn bộ quá trình này về bản chất vẫn tuân thủ chặt chẽ chu trình FDE:

\begin{center}
    \begin{minipage}{0.9\textwidth}
        \begin{enumerate}
            \item \textbf{[Fetch - Pha 1 đến 4]:} Giai đoạn nạp lệnh, CPU xuất địa chỉ từ PC ra bộ nhớ, 
            đọc mã lệnh 32-bit và đưa lệnh vào IR.
            \item \textbf{[Decode \& Update - Pha 5]:} Giải mã lệnh, Controller phân tích 3-bit opcode để phát tín hiệu điều khiển. 
            Tại thời điểm này, PC tự động tăng lên 1 ($PC \leftarrow PC + 1$) để trỏ sẵn tới lệnh kế tiếp.
            \item \textbf{[Execute - Pha 6 và 7]:} Thực thi, ALU lấy toán hạng từ bộ nhớ (nếu có) và tiến hành tính toán. 
            Đối với lệnh rẽ nhánh, PC sẽ được cập nhật lại địa chỉ nhảy nếu thỏa mãn điều kiện cờ (ví dụ lệnh SKZ).
            \item \textbf{[Write-back - Pha 8]:} Lưu trữ, kết quả tính toán được đưa vào Accumulator, 
            hoặc dữ liệu từ Accumulator được ghi vào bộ nhớ (đối với lệnh STO).
        \end{enumerate}
    \end{minipage}
\end{center}

\subsection{Tập lệnh và mã hóa lệnh}
CPU sử dụng định dạng lệnh 8-bit, trong đó 3-bit cao là opcode và 5-bit thấp là toán hạng (địa chỉ):

\begin{table}[H]
\centering
\begin{tabular}{|c|c|c|c|c|c|c|c|}
\hline
7 & 6 & 5 & 4 & 3 & 2 & 1 & 0 \\ \hline
\multicolumn{3}{|c|}{\textbf{Opcode (3 bits)}} & \multicolumn{5}{c|}{\textbf{Operand (5 bits)}} \\ \hline
\end{tabular}
\caption{Định dạng mã hóa lệnh 8-bit}
\end{table}

\begin{table}[H]
\centering
\renewcommand{\arraystretch}{1.2}
\small
\setlength{\tabcolsep}{3pt}
\begin{tabular}{|l|c|l|l|}
\hline
\textbf{Opcode} & \textbf{Mã} & \textbf{Hoạt động} & \textbf{Kết quả (Output)} \\ \hline
\textbf{HLT} & 000 & Dừng hoạt động chương trình & inA \\ \hline
\textbf{SKZ} & 001 & Nếu ALU = 0, bỏ qua lệnh tiếp theo & inA \\ \hline
\textbf{ADD} & 010 & Cộng Accumulator với bộ nhớ & inA + inB \\ \hline
\textbf{AND} & 011 & AND Accumulator với bộ nhớ & inA AND inB \\ \hline
\textbf{XOR} & 100 & XOR Accumulator với bộ nhớ & inA XOR inB \\ \hline
\textbf{LDA} & 101 & Đọc giá trị từ địa chỉ vào Accumulator & inB \\ \hline
\textbf{STO} & 110 & Ghi Accumulator vào địa chỉ bộ nhớ & inA \\ \hline
\textbf{JMP} & 111 & Nhảy không điều kiện đến địa chỉ đích & inA \\ \hline
\end{tabular}
\caption{Tập lệnh của RISC CPU}
\end{table}

\subsection{Ứng dụng thực tế của kiến trúc RISC}
Hiện nay kiến trúc RISC được ứng dụng rộng rãi trong thiết kế hệ thống số, từ các vi điều khiển nhúng đơn giản đến các bộ vi xử lý hiệu năng cao. 
Sự tối ưu về năng lượng tiêu thụ và diện tích phần cứng là yếu tố chính giúp RISC trở thành tiêu chuẩn phổ biến trong công nghiệp.
\begin{itemize}
    \item \textbf{Hệ sinh thái ARM (Advanced RISC Machine):}
    \begin{itemize}
        \item Thiết bị di động: Chiếm hơn 95\% thị phần smartphone toàn cầu (Qualcomm Snapdragon, Apple A-series, MediaTek Dimensity).
        \item Máy tính cá nhân: Bước ngoặt với Apple Silicon (M1, M2, M3 series), chứng minh RISC có thể vượt mặt CISC (x86) về hiệu năng trên mỗi Watt.
        \item Điện toán đám mây: Các bộ vi xử lý như AWS Graviton hay Ampere Altra đang dần thay thế CPU truyền thống trong các trung tâm dữ liệu.
    \end{itemize}

    \item \textbf{Kiến trúc MIPS (Microprocessor without Interlocked Pipeline Stages):}
    \begin{itemize}
        \item Thiết bị mạng: Phổ biến trong các bộ định tuyến (Router), thiết bị chuyển mạch (Switch) của Cisco, Linksys và TP-Link.
        \item Hệ thống nhúng: Sử dụng trong các bộ điều khiển công nghiệp, máy chơi game cổ điển (PS1, N64) và các thiết bị IoT yêu cầu độ ổn định cao.
        \item Giáo dục: Là mô hình chuẩn để giảng dạy môn Kiến trúc máy tính nhờ cấu trúc pipeline 5 giai đoạn kinh điển.
    \end{itemize}

    \item \textbf{RISC-V (Kiến trúc tập lệnh mở):}
    \begin{itemize}
        \item Tùy biến chuyên dụng: Các công ty như Western Digital và NVIDIA sử dụng RISC-V để thiết kế các bộ điều khiển (controller) tùy chỉnh cho ổ cứng và card đồ họa.
        \item Trí tuệ nhân tạo (AI): RISC-V cho phép thêm các tập lệnh tăng tốc AI riêng biệt, giúp tối ưu hóa việc xử lý dữ liệu lớn.
        \item Chủ quyền công nghệ: Các quốc gia và tập đoàn lớn đang đầu tư vào RISC-V để giảm sự phụ thuộc vào các chuẩn đóng như ARM hay x86.
    \end{itemize}
\end{itemize}
